% ===================================================================
%  TABELL Nr. 11 — D'Stad an hir Monumenter. --- De Brand.
%  Traduction 2026 d'apres texto/io, controlee sur texto/fr, texto/de
%  et texto/nl. Consigne : voir texto/lb/10-tabell-01.tex.
%
%  DEUX INSTITUTIONS DE LA MEME PAGE, ET DEUX SOLUTIONS OPPOSEES.
%  Les colonnes tcheque et lituanienne avaient trouve, au meme
%  tableau, deux mots refaits contre le voisin : « spořitelna » et
%  « hasiči », « taupomoji kasa » et « ugniagesiai » — la caisse
%  d'epargne et les eteigneurs. LE LUXEMBOURGEOIS EN FAIT UN ET
%  EMPRUNTE L'AUTRE.
%
%  La caisse d'epargne est « Spuerkeess », compose luxembourgeois de
%  « spueren » epargner et « Keess » caisse. Les pompiers sont
%  « Pompjeeën », c'est-a-dire les pompiers francais, sans un
%  changement au-dela de la finale.
%
%  ET CE N'EST PAS UNE INCOHERENCE, C'EST LA LOI DU TABLEAU 10 PRISE
%  SUR DES INSTITUTIONS. La Sparkasse est venue d'Allemagne avec un
%  nom allemand, que le luxembourgeois refait sans effort parce que
%  les deux morceaux existaient chez lui. Les pompiers sont venus de
%  France en corps constitue, avec leur nom : il n'y avait rien a
%  refaire, et personne n'a essaye. CE QUI ARRIVE TOUT MONTE ARRIVE
%  AVEC SON NOM.
%
%  ET LE RESTE DE L'ADMINISTRATION SUIT LA MEME REGLE SANS EXCEPTION.
%  « Gemeng » la commune, « Gemengenhaus » la mairie,
%  « Buergermeeschter » le bourgmestre, « Geriicht » le tribunal,
%  « Spidol » l'hopital sont germaniques ; « Bourse »,
%  « Distriktskommissariat », « Policekommissär », « Ambulanz »,
%  « Affiche », « Peristyl », « Café », « Kiosk » sont francais. On
%  peut lire dans le nom d'un bureau de quel cote de la frontiere il
%  est entre.
% ===================================================================

%%K t11-apar-1 apar
\VUsaut{2.0mm}
\VUcentre{13.2pt}{90}{DRËTT SERIE}
\VUsaut{3.0mm}
\VUfilet{16mm}
\VUsaut{5.6mm}
\VUtitre{15.0pt}{60}{TABELL Nr. 11}
\VUsaut{1.4mm}
\VUpk{11.4pt}{40}{D'Stad an hir Monumenter. --- De Brand.}
\VUsaut{2.2mm}

%%K t11-tit-1 sub
\VUpk{10.2pt}{40}{D'Virstad.}
\VUsaut{3.0mm}

%%K t11-c1-01-1 p
1. --- Ech wunnen an enger grousser Stad, déi 100 Kilometer vum Ozean ewech läit an
déi trotzdeem en éischtrangegen \VUgras{Hafen} \textsuperscript{(1)} ass. De Floss,
deen se säumt, ass 400 Meter breet a bilt eng ganz schéi Rieds.

%%K t11-c1-02-1 p
2. --- Deen ganzen Deel vun der Stad, deen op de \VUgras{Kaien} \textsuperscript{(2)}
läit, gesäit ganz maritim an industriell aus a bilt eng \VUgras{Virstad}
\textsuperscript{(3)}, déi nëmme vu Schëffsleit an Aarbechter bewunnt gëtt. Dës
schaffen an de \VUgras{Fabriken} \textsuperscript{(4)} an am \VUgras{Gaswierk}
\textsuperscript{(5)}, vun deem de héije \VUgras{Kamäin} \textsuperscript{(6)} an de
\VUgras{Gasometer} vu ganz wäitem ze gesi sinn.

%%K t11-c1-03-1 p
3. --- Am Mëttelalter war d'Stad befestegt, an ee fënnt nach e puer Reschter vun
hire Wäll, virun allem d'\VUgras{Paart} \textsuperscript{(8)} vum alen
\VUgras{befestegte Schlass} \textsuperscript{(7)}, vun hiren zwee \VUgras{Tierm}
\textsuperscript{(9)} flankéiert, déi elo als \VUgras{Prisong} benotzt ginn.

%%K t11-c1-04-1 p
4. --- An deem ganze \VUgras{Quartier,} dat ganz al ausgesäit, ass nëmmen ee Gebai
relativ modern: dat ass d'\VUgras{Zollamt} \textsuperscript{(10)}. Et läit tëscht dem
Kai an der \VUgras{Gaass} \textsuperscript{(11)}, déi bei dat aalt
\VUgras{Gemengenhaus} \textsuperscript{(12)} féiert. Ee erkennt dee leschte Monument
liicht um risege \VUgras{Klackentuerm} \textsuperscript{(13)}, deen déi al Stad
beherrscht.

%%K t11-c1-05-1 p
5. --- Op der anerer Säit vun der Strooss \VUgras{steet} e Gebai, dat am selwechte
Stil wéi d'Zollamt gebaut ass: dat ass d'\VUgras{Bourse} \textsuperscript{(14)}. Eng
vu senge Fassade kuckt op eng kleng Plaz, wou d'\VUgras{Distriktskommissariat}
\textsuperscript{(15)} ass. Tatsächlech ass eis Stad, ouni eng Haaptstad ze sinn,
trotzdeem e wichtegen Distriktshaaptuert.

%%K t11-tit-2 sub
\VUsaut{5.0mm}
\VUpk{10.2pt}{40}{Deen héchsten Deel vun der Stad.}
\VUsaut{3.4mm}

%%K t11-c1-06-1 p
6. --- D'Garnisoun, zimlech zahlräich, ass a grousse ganz gesonde \VUgras{Kasären}
\textsuperscript{(16)} ënnerbruecht; si zéie sech bis un d'Gitter vum
\VUgras{Gemengepark} \textsuperscript{(17)}. De \VUgras{Boulevard}
\textsuperscript{(19)}, deen an dee Park féiert, geet hannert der
\VUgras{Spuerkeess} \textsuperscript{(18)} laanscht a mündet op der Plaz vun der
\VUgras{Kathedral} \textsuperscript{(20)}.

%%K t11-c1-07-1 p
7. --- All déi Monumenter zéie sech amphitheatralesch op engem Hiwwelchen, wou och
eise groussen \VUgras{Lycée} \textsuperscript{(21)} ass.

%%K t11-c1-07-2 p
Wann een e liicht ofschësseg Wee erofgeet, deen op déi Grouss Gaass kënnt, kënnt een
bei d'\VUgras{Geriicht} \textsuperscript{(22)}, virun deem sech en agreabelen
\VUgras{ëffentleche Gaart} \textsuperscript{(26)} zitt. Dee \VUgras{Gaardeplaz} ass
net ganz grouss, mä an der Mëtt vun der Stad mécht en eng Oas vu Gréngs a Killt.
Dofir gëtt en immens vill vun de Stadleit besicht, déi gär ënner dem Zelt vum
\VUgras{Café} \textsuperscript{(25)} setze kommen, fir de Militärmuseker
nozelauschteren, déi donneschdes a sonndes am \VUgras{Kiosk} \textsuperscript{(27)}
spillen, deen tëscht de Wisen an de Blummebeeter steet.

%%K t11-c1-08-1 p
8. --- Op enger vun deene Wise steet eng bronzen \VUgras{Statu} \textsuperscript{(28)},
e Monument, dat zur Erënnerung un ee vun eise Matbierger opgeriicht gouf, deen
senger Gebuertsstad grouss Déngschter geleescht huet.

%%K t11-tit-3 sub
\VUsaut{4.0mm}
\VUpk{10.2pt}{40}{D'Transportmëttelen.}
\VUsaut{3.4mm}

%%K t11-c1-09-1 p
9. --- Bis elo ass eis Stad nach net mat elektreschem Liicht beliicht, si huet nëmme
\VUgras{Gaslampen} \textsuperscript{(29)}. Mä d'\VUgras{Zeitunge vun hei} maachen eng
Kampagne, fir datt dat Beliichtungssystem verbessert gëtt, \VUgras{wéi} een schonn
\VUgras{d'Traktiounssystem vun den Tramen} verbessert huet. \VUgras{Déi} al Wagonen,
\VUgras{gezunn} vu \VUgras{Päerd,} goufen praktesch duerch \VUgras{elektresch Tramen}
\textsuperscript{(31)} ersat, déi d'Reesend séier vun engem Enn vun der
\VUgras{Stad bis bei dat anert} transportéieren, fir e ganz \VUgras{bëllege}
Präis.

%%K t11-c1-10-1 p
10. --- Nieft dem Geriichtsgebai ass d'\VUgras{Bank} \textsuperscript{(32)}, wou
d'Handelspabeieren diskontéiert ginn. D'\VUgras{Bankboten} \textsuperscript{(33)}
ginn d'Scholde bei de Leit doheem anzéien.

%%K t11-tit-4 sub
\VUsaut{4.0mm}
\VUpk{10.2pt}{40}{Konscht a Bildung.}
\VUsaut{3.4mm}

%%K t11-c1-11-1 p
11. --- Dat schéint Gebai, dat nieft drun steet, ass de \VUgras{Musée.} Et huet
zugläich d'\VUgras{Bibliothéik} \textsuperscript{(34)} vun der Stad, déi schéin
\VUgras{Naturwëssenschaftssammlungen} \textsuperscript{(35)} (Naturmusée) an e
zierleche \VUgras{Biller-Musée} \textsuperscript{(36)}, deen eng herrlech stänneg
\VUgras{Ausstellung} bilt.

%%K t11-c1-11-2 p
En ass zweemol d'Woch fir de Publikum op, mä d'Auslänner dierfen en all Dag besichen,
géint e klengt Trinkgeld fir de \VUgras{Wiechter} \textsuperscript{(37)}. D'\VUgras{Moler}
\textsuperscript{(38)} kommen do schaffen. Ee gesäit se \VUgras{virun} hire
Staffelaien, mat der \VUgras{Palette} an der Hand, wéi si déi berühmt Biller
kopéieren.

%%K t11-c1-12-1 p
12. --- Op der Héicht, hannert dem Musée, gesäit ee d'\VUgras{Kuppel}
\textsuperscript{(40)} vum \VUgras{Spidol} \textsuperscript{(39)} an d'Gebaier vun der
\VUgras{Normalschoul} \textsuperscript{(41)}, an där d'Léierpersonal vun eise
Gemengeschoulen ausgebilt gouf.

Op der Theaterplaz ass e \VUgras{Postbüro} \textsuperscript{(43)}, dee an engem
\VUgras{Privathaus} \textsuperscript{(42)} agerecht ass.

%%K t11-tit-5 sub
\VUsaut{5.0mm}
\VUpk{11.4pt}{40}{De Brand.}
\VUsaut{3.0mm}

%%K t11-tit-6 sub
\VUpk{10.2pt}{40}{Feiergejäits.}
\VUsaut{3.4mm}

%%K t11-c2-01-1 p
1. --- Eng schrecklech Neiegkeet verbreet sech duerch d'Stad: am Theater ass grad e
Brand ausgebrach! Am Ablack, wou ech op der \VUgras{Plaz} \textsuperscript{(96)}
ukommen, ass d'Masse schonn um \VUgras{Trottoir} \textsuperscript{(46)} an op der
\VUgras{Chaussée} \textsuperscript{(47)} versammelt. D'\VUgras{Postbeamten}
\textsuperscript{(44)} an d'\VUgras{Bréifdréier} \textsuperscript{(45)} kommen aus
dem Postbüro. E \VUgras{Tramkutscher} \textsuperscript{(48)} huet misse säi Wagon
uhalen, fir den \VUgras{Ambulanzwon} \textsuperscript{(50)} vun der Stad
laanschtzeloossen, deen mat engem \VUgras{Chirurg} \textsuperscript{(52)} an engem
\VUgras{Infirmier} \textsuperscript{(51)} ukënnt, fir e \VUgras{Blesséierten}
\textsuperscript{(53)} ze verspriechen, deen op enger \VUgras{Bahre}
\textsuperscript{(54)} nieft de \VUgras{Zeitungskiosk} \textsuperscript{(55)} bruecht
gouf.

%%K t11-c2-02-1 p
2. --- Eng Infanteriekompanie, déi vun den Übungen zréckkoum, huet ugehalen, fir mat
de \VUgras{beridden Stadgendaarmen} \textsuperscript{(61)} den Uerdnungsdéngscht ze
sécheren. \VUgras{D'Reider,} deenen hir Päerd sech opriichten, dreiwe besser wéi
d'\VUgras{Zaldoten} \textsuperscript{(57)} oder d'\VUgras{Gendaarmen}
\textsuperscript{(63)} d'\VUgras{Neigiereg} \textsuperscript{(56)} zréck. Iwwregens sinn
d'Gendaarme ganz
beschäftegt. Ee vun hinne weist de \VUgras{Polizisten} \textsuperscript{(64)}, wouhinner
een en anere \VUgras{Verletzten} \textsuperscript{(65)} bréngen muss, e mutege
Pompjee, dee vun enger Leeder gefall ass.

%%K t11-c2-03-1 p
3. --- Den \VUgras{Offizéier} \textsuperscript{(66)} gëtt genau d'Plaz un, wou een
d'\VUgras{Dampspritz} \textsuperscript{(67)} opstelle muss, déi ukënnt. Während hien
op déi mächteg Maschinn waart, huet hien an aller Eil eng \VUgras{Handspritz}
\textsuperscript{(68)} opstelle gelooss, vun där de laange \VUgras{Schlauch}
\textsuperscript{(69)} d'Waasser a Strëmm op d'Feier geheit.

Sechs Pompjeeë bedéngen se, während hire Kamerad, deen d'\VUgras{Spritz}
\textsuperscript{(70)} hält, de \VUgras{Strahl} \textsuperscript{(71)} an d'Mëtt vum
Brand riicht.

%%K t11-c2-04-1 p
4. --- Op der anerer Säit vum Theater huet ee déi grouss \VUgras{Leeder}
\textsuperscript{(72)} ugeleet. Si mécht et de Pompjeeë méiglech, bei d'Flamen ze
kommen, déi tëscht Wirbele vu schwaarzem \VUgras{Rauch} \textsuperscript{(75)}
erausspritzen.

%%K t11-tit-7 sub
\VUsaut{5.0mm}
\VUpk{10.2pt}{40}{Mutteg Doten.}
\VUsaut{3.4mm}

%%K t11-c2-05-1 p
5. --- De \VUgras{Brand} \textsuperscript{(74)} ass am Foyer vum \VUgras{Theater}
\textsuperscript{(73)} entstanen, während ee d'\VUgras{Oper} geprouft huet, vun där
déi éischt \VUgras{Virstellung} muer sollt sinn. Glécklecherweis waren nëmme wéineg
\VUgras{Zuschauer} \textsuperscript{(76)} am Zuschauersall. Déi Aarm hu sech op de
\VUgras{Balkon} \textsuperscript{(77)} geflücht, wou mutteg \VUgras{Pompjeeën}
\textsuperscript{(86)} hinne mat Leederen a Seeler ze Hëllef kommen.

%%K t11-c2-06-1 p
6. --- Ënner dem \VUgras{Peristyl} \textsuperscript{(78)}, wou d'\VUgras{Affiche}
\textsuperscript{(79)} iwwert d'Virstellung informéiert, gesäit ee
d'\VUgras{Museker} \textsuperscript{(81)} vum \VUgras{Orchester}
\textsuperscript{(80)} flüchten, andeems si hir Instrumenter matdroen: eng
\VUgras{Gei} \textsuperscript{(81)}, eng \VUgras{Flütt} \textsuperscript{(82)}, e
\VUgras{Cello} \textsuperscript{(83)}, eng \VUgras{Posaun} \textsuperscript{(84)},
eng \VUgras{Trommel} \textsuperscript{(85)} asw. Ee versicht, en Deel vun de
\VUgras{Miwwelen} \textsuperscript{(87)} ze retten.

%%K t11-c2-07-1 p
7. --- D'\VUgras{Schauspiller} \textsuperscript{(89)} an d'\VUgras{Schauspillerinnen}
\textsuperscript{(88)}, d'\VUgras{Sänger} an d'\VUgras{Sängerinnen} hu misse mat
hirem Theater-\VUgras{Kostüm} \textsuperscript{(90)} flüchten. Den Tenor erkläert engem
\VUgras{Journalist} \textsuperscript{(92)}, wéi et geschitt ass. De \VUgras{Reporter}
ass vum éischten Ablack un mat den Autoritéiten erbäigelaf: dem
\VUgras{Distriktskommissär} \textsuperscript{(91)}, dem \VUgras{Buergermeeschter}
\textsuperscript{(93)}, dem \VUgras{Policekommissär} \textsuperscript{(94)} an dem
\VUgras{Untersuchungsriichter} \textsuperscript{(95)}, déi eng Enquête maache
wäerten, fir iwwert d'Verantwortlechkeeten ze urteelen.

\VUsaut{6.0mm}
\VUfilet{22mm}
